\subsection{Traditional Light Field Depth Estimation Methods}

Many depth estimation methods for light-field cameras have been proposed, primarily leveraging the geometric properties of epipolar plane images (EPIs) \cite{bolles1987epipolar, wanner2012globally}. These traditional approaches can be broadly categorized into EPI-based geometric methods and multi-cue integration techniques.

\subsubsection{EPI-Based Geometric Methods}
EPI-based approaches demonstrate remarkable performance under ideal imaging conditions by exploiting the strict photo-consistency across angular views. Wanner and Goldluecke \cite{wanner2012globally} proposed a globally consistent depth labeling framework by analyzing the dominant orientation of EPI structures using structure tensors. Specifically, the slope of the linear structures in an EPI is explicitly inversely proportional to the scene depth. The depth $d$ at a given pixel can be derived from the EPI slope $m$ as:
\begin{equation}
d = \frac{f \cdot B}{m \cdot \Delta u},
\label{eq:ch2_1}
\end{equation}
where $f$ is the focal length, $B$ is the baseline, and $\Delta u$ is the angular sampling interval. Building upon this geometric foundation, Kim et al.\ \cite{kim2013adaptive} introduced an adaptive thresholding technique to refine the slope estimation, successfully improving depth accuracy in textureless regions. 

Despite their success, these methods fundamentally rely on the Lambertian assumption. As analyzed in our work, EPI methods fail on non-Lambertian surfaces because specular and scattering reflections break the Lambertian angular cone geometric constraints. Instead of forming continuous linear structures, non-Lambertian surfaces generate X-shaped patterns or disrupted trajectories in the EPI space. This geometric root cause leads to severe depth estimation degradation, confirming that the failure stems from assumption violations rather than parameter defects.

\subsubsection{Multi-Cue Integration Methods}
While EPI-based methods predominantly focus on angular correspondence, Tao et al.\ \cite{tao2013depth} extended the cue integration by combining defocus and correspondence cues using light-field cameras. In their approach, the depth is derived by minimizing a joint energy function that aggregates both the defocus cue (measuring the sharpness of refocused images) and the correspondence cue (evaluating pixel similarity across sub-aperture images). The joint energy function $E(d)$ is formulated as:
\begin{equation}
E(d) = \sum_{x} \left( \lambda_{def} E_{def}(x, d) + \lambda_{cor} E_{cor}(x, d) \right) + E_{smooth}(d),
\label{eq:ch2_2}
\end{equation}
where $E_{def}$ and $E_{cor}$ represent the defocus and correspondence energy terms, respectively, $\lambda_{def}$ and $\lambda_{cor}$ are weighting parameters, and $E_{smooth}$ enforces spatial smoothness. 

Although integrating defocus cues provides supplementary information in textureless regions where correspondence fails, the defocus cue itself is highly sensitive to the material's reflectance properties. For non-Lambertian surfaces, the point spread function varies significantly with viewing angles, rendering the defocus sharpness metrics unreliable.

\subsubsection{Summary and Limitations}
To summarize the characteristics of traditional methods, we categorize their strengths and limitations in Table \ref{tab:traditional_methods}. The primary limitations of traditional methods can be summarized as follows:
\begin{itemize}
\item \textbf{Lambertian Assumption:} Implicit reliance on diffuse reflectance, causing failures on specular and scattering surfaces.
\item \textbf{Geometric Breakdown:} Inability to handle X-shaped EPI patterns generated by non-Lambertian materials.
\item \textbf{Cue Ambiguity:} Defocus and correspondence cues become unreliable under complex light transport.
\end{itemize}

This physical limitation boundary indicates that merely tuning parameters or loss functions is insufficient; instead, a fundamental architectural shift that explicitly models non-Lambertian physical properties is required, which motivates our proposed dual-branch routing mechanism.

\begin{table}[!t]
\centering
\caption{Comparison of Traditional Light Field Depth Estimation Methods}
\label{tab:traditional_methods}
\begin{tabular}{@{}llll@{}}
\toprule
\textbf{Category} & \textbf{Method} & \textbf{Strengths} & \textbf{Limitations} \\ \midrule
EPI Geometry & Structure Tensor \cite{wanner2012globally} & High accuracy in textured regions & Fails on non-Lambertian X-shaped EPIs \\
Multi-Cue & Defocus \& Correspondence \cite{tao2013depth} & Robust in textureless areas & Defocus unreliable for specular surfaces \\ \bottomrule
\end{tabular}
\end{table}

\subsection{Deep Learning-based Methods}

With the advent of deep learning, convolutional neural networks (CNNs) and vision transformers have significantly advanced light field depth estimation by learning complex mappings from high-dimensional light field data to depth maps. These methods can be broadly categorized based on their input representations and architectural designs.

\subsubsection{EPI-based Deep Networks}
Epipolar plane images (EPIs) remain a dominant representation in deep learning approaches due to their explicit encoding of geometric disparities. Shin et al.\ \cite{shin2018epinet} proposed EPINet, which extracts 4-directional EPIs and employs a multi-stream CNN to regress depth. The network learns to identify the linear slopes in EPIs, effectively replacing traditional structure tensor calculations with data-driven feature extraction. Subsequent works, such as LFNet \cite{liang2020lfnet} and AGL-Net \cite{wang2021agl}, introduced attention mechanisms and adaptive geometric learning to enhance the extraction of EPI features, particularly in occluded and textureless regions.

Despite their success on Lambertian datasets, EPI-based deep networks inherit the physical limitations of the EPI representation itself. As established in our geometric root cause analysis, non-Lambertian surfaces (e.g., specular and scattering materials) disrupt the photo-consistency constraint, generating X-shaped or curved trajectories in the EPI space rather than straight lines. Consequently, standard EPI-based CNNs struggle to extract valid disparity cues from these anomalous patterns, leading to severe depth degradation in non-Lambertian regions.

\subsubsection{Macro-Pixel and Multi-View Methods}
To circumvent the high computational cost and information loss associated with EPI extraction, several methods directly process macro-pixels (sub-aperture images arranged in a 2D grid) or multi-view stacks. Wang et al.\ \cite{wang2018occlusion} proposed a macro-pixel image (MacPI) based network that applies 4D convolutions to capture both spatial and angular correlations simultaneously. Similarly, methods adapting stereo matching architectures, such as PSMNet \cite{chang2018pyramid} and GANet \cite{zhang2019ga}, construct 3D cost volumes by warping sub-aperture images across hypothetical depth planes. The cost volume $C(x, y, d)$ is typically aggregated using:
\begin{equation}
C(x, y, d) = \frac{1}{|\mathcal{V}|} \sum_{(u,v) \in \mathcal{V}} \phi(I_{u,v}(x + u \cdot d, y + v \cdot d)),
\label{eq:ch2_3}
\end{equation}
where $\mathcal{V}$ denotes the set of angular views, $d$ is the hypothesized disparity, and $\phi$ represents a feature extraction function.

While cost volume-based methods excel at capturing global context and handling occlusions, they implicitly assume Lambertian reflectance during the warping process. For non-Lambertian surfaces, the view-dependent appearance changes violate the photo-consistency assumption, resulting in blurred or multi-modal cost volume distributions. This ambiguity severely hinders the subsequent 3D convolution or regularization steps from identifying the correct depth peak.

\subsubsection{Attention and Transformer-based Architectures}
Recently, vision transformers and self-attention mechanisms have been introduced to light field depth estimation to model long-range dependencies across angular views. LF-Transformer \cite{liang2022light} and Epinet-Transformer \cite{chen2023epinet} utilize self-attention to dynamically weight the contributions of different sub-aperture images, effectively suppressing occluded or inconsistent views. The attention weights $\alpha_{i,j}$ between spatial-angular tokens are computed as:
\begin{equation}
\alpha_{i,j} = \text{softmax}\left(\frac{Q_i K_j^T}{\sqrt{d_k}}\right),
\label{eq:ch2_4}
\end{equation}
where $Q_i$ and $K_j$ are the query and key projections of the $i$-th and $j$-th tokens, respectively.

Although attention mechanisms improve robustness against occlusions and view inconsistencies, they treat non-Lambertian anomalies merely as noise or outliers to be down-weighted, rather than modeling their underlying physical properties. Furthermore, the quadratic computational complexity of self-attention with respect to the number of views restricts these models from fully exploiting high-angular-resolution light fields (e.g., $9 \times 9$ or higher), often forcing them to downsample the angular dimensions and lose critical high-frequency material cues.

\subsubsection{Physics-Informed and Material-Aware Networks}
Recognizing the limitations of purely data-driven approaches, recent studies have begun to incorporate physical priors into network design. Some works integrate defocus cues or shading models as auxiliary tasks to guide depth regression \cite{huang2021physics}. However, these methods typically rely on heuristic fusion strategies and lack a unified physical model to explicitly decouple illumination, geometry, and material reflectance. 

To address the non-Lambertian challenge, a few dual-branch or multi-branch architectures have been explored to separate diffuse and specular components \cite{zhang2022specular}. Nevertheless, these approaches often depend on explicit specular mask annotations, which are scarce and difficult to obtain in real-world light field datasets. They also fail to provide a theoretical foundation for feature routing, relying instead on end-to-end learning to implicitly discover the boundary between Lambertian and non-Lambertian regions.

\subsubsection{Summary and Limitations}
Table \ref{tab:dl_methods} summarizes the characteristics of existing deep learning-based methods. The primary limitations can be concluded as follows:
\begin{itemize}
\item \textbf{Implicit Lambertian Bias:} Most architectures implicitly encode the Lambertian assumption in their feature extraction or cost volume construction, lacking explicit mechanisms to handle non-Lambertian physical anomalies.
\item \textbf{Lack of Physical Decoupling:} Existing physics-informed methods fail to rigorously decouple the angular signal into illumination, geometry, and material components, leading to suboptimal feature representations.
\item \textbf{Ineffective Routing Mechanisms:} Current multi-branch designs rely on heuristic or data-driven routing without geometric pseudo-labels, resulting in unstable optimization and poor generalization across diverse material domains.
\end{itemize}

These limitations highlight the critical need for a unified physical model that explicitly decomposes the angular signal and employs a geometry-aware dual-mask routing mechanism, which forms the core contribution of our proposed GeometricDualMask architecture.

\begin{table*}[!t]
\centering
\caption{Comparison of Deep Learning-based Light Field Depth Estimation Methods}
\label{tab:dl_methods}
\resizebox{\textwidth}{!}{
\begin{tabular}{@{}lllll@{}}
\toprule
\textbf{Category} & \textbf{Representative Methods} & \textbf{Input Representation} & \textbf{Key Innovations} & \textbf{Limitations on Non-Lambertian Surfaces} \\ \midrule
EPI-based & EPINet \cite{shin2018epinet}, LFNet \cite{liang2020lfnet} & 4-directional EPIs & Data-driven slope extraction, attention & Fails on X-shaped EPI patterns; assumes linear trajectories \\
Macro-Pixel / Cost Volume & MacPI \cite{wang2018occlusion}, PSMNet \cite{chang2018pyramid} & Sub-aperture images, 3D cost volume & 4D convolutions, global context aggregation & Photo-consistency violation blurs cost volume peaks \\
Transformer-based & LF-Transformer \cite{liang2022light} & Spatial-angular tokens & Long-range angular dependency modeling & High computational cost; treats material anomalies as noise \\
Physics-Informed & Defocus-guided \cite{huang2021physics}, Dual-branch \cite{zhang2022specular} & Multi-cue, multi-branch & Auxiliary physical cues, component separation & Lacks unified angular signal decomposition; requires rare masks \\ \bottomrule
\end{tabular}
}
\end{table*}

\subsection{Summary and Motivation}

Despite the remarkable progress achieved by existing light field depth estimation methods, a critical gap remains in handling complex real-world scenes characterized by non-Lambertian surfaces. As summarized in Section 2.1 and Section 2.2, current approaches predominantly rely on the Lambertian assumption, which enforces strict photo-consistency and linear trajectories in Epipolar Plane Images (EPIs). However, this assumption is fundamentally violated in the presence of specular reflections, glossy materials, and complex scattering, leading to severe performance degradation. 

Through an exhaustive exploration of 132 experiments across 16 research directions, we identify that the performance bottleneck on non-Lambertian surfaces is rooted in \textit{assumption violations} rather than mere \textit{parameter defects} (e.g., suboptimal learning rates or loss weights). Specifically, we pinpoint three fundamental gaps in the state-of-the-art that motivate our work:

\begin{itemize}
    \item \textbf{Inadequate Physical Decoupling under Discrete Angular Sampling:} Existing physics-informed methods often attempt to separate diffuse and specular components using frequency-domain analyses (e.g., 2D-DFT). However, these spectral methods fail under the low-resolution discrete angular sampling (e.g., $9 \times 9$ views) typical of commercial light field cameras. There is a lack of a robust physical model that can decouple the angular signal into illumination, geometry, and material components directly in the spatial-angular domain.
    \item \textbf{Misunderstanding of EPI Geometric Degradation:} Traditional EPI-based networks treat non-Lambertian anomalies as noise or occlusions. Our analysis reveals that specular and scattering surfaces destroy the Lambertian angular cone geometric constraints, forming distinct X-shaped patterns rather than linear structures in EPIs. Current architectures lack a geometry-aware routing mechanism to explicitly divert these X-shaped patterns to specialized processing branches, resulting in unstable optimization.
    \item \textbf{Scarcity of Cross-Domain Generalization Frameworks:} Most deep learning models are trained on homogeneous datasets (e.g., purely Lambertian synthetic scenes) and suffer from severe overfitting when exposed to non-Lambertian data. There is an urgent need for a unified dataset loader and a domain-balanced training strategy that can statistically validate the extracted physical features and ensure robust generalization across diverse material domains.
\end{itemize}

To bridge these gaps, we propose a \textbf{Unified Dual-Mask Physical Model} for non-Lambertian light field depth estimation. Our motivation is threefold. First, we introduce a physics-based three-layer angular signal decomposition model, formulated as:
\begin{equation}
    I(u,v) = \text{DC} + a \cdot u + b \cdot v + \varepsilon(u,v),
    \label{eq:ch2_1}
\end{equation}
where the angular signal $I(u,v)$ is decoupled into ambient illumination (DC term), disparity/depth (linear terms $a$ and $b$), and material reflectance (residual term $\varepsilon(u,v)$). By extracting geometric angular features such as the angular gradient from the residual term, we establish a reliable foundation for material classification without relying on ineffective spectral analysis.

Second, motivated by the geometric root cause of EPI failure, we design a dual-branch adaptive routing mechanism. Instead of heuristic routing, we utilize the extracted geometric pseudo-labels to generate pixel-wise dual masks, effectively routing Lambertian regions to an EPI-based branch and non-Lambertian regions (characterized by X-shaped patterns and bimodal distributions) to a dedicated geometric branch. 

Finally, to ensure cross-domain robustness, we construct a unified light field dataset encompassing approximately 230 scenes across five major domains (Urban, Mixed, Lambertian, Non-Lambertian, etc.). We perform global statistical validation on the decomposed features using Cohen's $d$ effect size and implement a multi-domain balanced sampling strategy to mitigate overfitting caused by the scarcity of non-Lambertian data. This comprehensive framework shifts the paradigm from blind parameter tuning to geometry-aware physical modeling, paving the way for unified and robust light field depth estimation.